\subsection{Problem Setup and Framework Overview} \label{sec:04-method-setup}
\begin{foldable} % Setup + overview merged
  TAKE instantiates the RDM objective (Eq.~\ref{eq:reweighted-distribution-matching}) by solving the two open problems from \S\ref{sec:03-theory-failures}.
  Let $\mathcal{D} = \{z_n\}_{n=1}^{N}$ denote the training corpus, $\mathcal{D}_\text{synth} = \{\tilde{z}_c\}_{c=1}^{M}$ the synthetic candidate pool, and $\tilde{\mathcal{D}} \subset \mathcal{D}_\text{synth}$ the distilled set of budget $K$.
  Indices $n \in [N]$ and $t \in [T]$ range over training samples and checkpoints respectively.
  TAKE proceeds in three stages:
  \begin{enumerate}
    \item \textbf{Score}---train a logistic probe on $\mathcal{D}$ across $T$ checkpoints, compute influence matrix $I \in \mathbb{R}^{N \times T}$, apply reciprocal reweighting and temporal weighting to obtain trajectory-aware knowledge scores $\kappa_n$.
    \item \textbf{Generate}---fine-tune a small language model on $\mathcal{D}$ and sample to produce the synthetic candidate pool $\mathcal{D}_\text{synth}$.
    \item \textbf{Distill}---solve discrete OT with source weights $\kappa_n$ over $\mathcal{D}_\text{synth}$ to obtain the distilled corpus $\tilde{\mathcal{D}}$ of budget $K$.
  \end{enumerate}
  The three stages are decoupled and independently replaceable, with $\phi$ shared across Stages 1 and 3 as a natural consistency choice. The following sections detail each in turn.
\end{foldable}

\subsection{Gradient-Based Influence Along the Trajectory} \label{sec:04-method-gradient}
\begin{foldable} % Motivation for trajectory-aware estimation
  The hard-sample bias (HSB; \S\ref{sec:03-theory-failures}) arises because single-checkpoint influence scores overemphasize noisy or difficult samples, whose gradients remain large at convergence.
  The fix is to integrate gradient information across the full training trajectory, producing a trajectory-aware knowledge score $\kappa_n$ per sample.
  We build $\kappa_n$ in two steps: constructing the influence matrix $I$ (\S\ref{sec:04-gradient-matrix}), then collapsing it into a per-sample scalar score via reciprocal reweighting and temporal weighting (\S\ref{sec:04-gradient-kappa}).
  The resulting $\kappa_n$ provides a task-aligned, bias-corrected weight for $P^w$.
\end{foldable}

\subsubsection{Influence Matrix} \label{sec:04-gradient-matrix}
\begin{foldable} % Definition and construction of I
  Our goal is to construct $I \in \mathbb{R}^{N \times T}$ whose entry $I_{(n,t)}$ reflects how much sample $z_n$ contributes to learning at checkpoint $\theta_t$.
  The influence self-score (\S\ref{sec:03-theory-failures}, Eq.~\ref{eq:influence-self}) is the natural candidate, but requires Hessian inversion at a converged checkpoint and is undefined mid-trajectory.
  We address both limitations by adapting the per-sample gradient norm from TracIn~\cite{pruthi2020estimating} to the self-influence setting, querying training points instead of test points:
  \begin{equation}
    I_{(n,t)} = \|\nabla_\theta \ell(z_n;\,\theta_t)\|.
    \label{eq:influence-matrix}
  \end{equation}
  As the $H=I$ special case of the influence self-score (Remark~\ref{rem:grad-norm}), $I_{(n,t)}$ exhibits the same HSB structure---the bias that $\kappa_n$ corrects via reciprocal inversion and trajectory integration (Proposition~\ref{prop:traj}).
  Scores are column-wise normalized $I_{(n,t)} \leftarrow I_{(n,t)} / \sum_{n'} I_{(n',t)}$ for comparability across checkpoints.

  In practice, we score samples with a linear head $W$ atop a pre-trained, frozen sentence encoder $\phi$, giving $\theta = (W, \phi)$.
  The signal quality comes not from the linear head but from $\phi(x_n)$: $\phi$ is pre-trained on large-scale corpora and $W$ is fine-tuned on the full task corpus, so the representations, and hence the gradient norms, are both semantically rich and task-aligned by construction.
  Following the standard practice in scalable influence estimation~\cite{guo2021fastif}, we approximate full-model influence via the last layer---here, the linear head $W$:
  \begin{equation}
    \nabla_\theta \ell(z_n;\,\theta_t) \approx \nabla_W \ell(z_n;\,\theta_t),
    \label{eq:last-layer}
  \end{equation}
  trading a small bias for tractability.
  By the Goodfellow outer-product identity~\cite{goodfellow2015efficient}, $\nabla_W \ell = g_n \phi(x_n)^\top$ where $g_n$ is the output gradient.
  The Goodfellow trick vectorizes the batch backward pass, avoiding per-sample loops.
  Adapting to a new task is trivial: retrain $W$ on the new corpus; $\phi$ and the influence pipeline are unchanged.
\end{foldable}

\subsubsection{Knowledge Score and Temporal Weighting} \label{sec:04-gradient-kappa}
\begin{foldable} % Inversion and temporal weighting to produce kappa
  To correct the HSB, we require a score that is high for well-fitted samples and low for noisy ones---the opposite of what $I_{(n,t)}$ directly provides.
  Since a well-fitted sample has a small gradient norm, taking the reciprocal naturally inverts this ranking.
  We further weight each checkpoint by a decreasing temporal kernel $k : \{0,\ldots,T-1\} \to \mathbb{R}_{>0}$, giving:
  \begin{equation}
    \kappa_n = \sum_{t=0}^{T-1} \frac{k(t)}{I_{(n,t)}},
    \label{eq:kappa}
  \end{equation}
\end{foldable}

\begin{foldable}
  A decreasing kernel is the principled choice: early checkpoints capture easy samples fitting rapidly, mid checkpoints offer the strongest discriminative signal, and late checkpoints see the gap narrow as memorization sets in~\cite{arpit2017closer}.
  Thus, concentrating mass on the early-to-mid phase and discounting late memorization directly corrects the HSB structure identified in Proposition~\ref{prop:hsb}.
  By default we use the exponential kernel $k(t) = e^{-\lambda t / T}$ ($\lambda > 0$), with $\lambda$ treated as a hyperparameter.
  Alternative kernels (linear, cosine) are drop-in replacements, and Proposition~\ref{prop:traj} guarantees bias reduction holds for any strictly decreasing $k$.
  The resulting $\kappa_n$ summarizes each sample's cumulative, bias-corrected learning contribution across the full trajectory, resolving open problem~\ref{open-problem-weight}.
  Normalized source weights $w_n = \kappa_n / \sum_{n'} \kappa_{n'}$ define the reweighted distribution $P^w$ and serve as input to the OT selection step (\S\ref{sec:04-method-ot}).
\end{foldable}

\begin{foldable}
  \begin{proposition}[Trajectory Integration Reduces HSB] \label{prop:traj}
    Under smooth loss and bounded gradient noise, $\kappa_n$ produces a strictly smaller noisy-over-clean gap than the single last-checkpoint baseline $\kappa^\dagger_n = 1/I_{(n,T-1)}$:
    \begin{equation}
      \mathbb{E}_{z \sim \mathcal{D}_\textup{noisy}}[\kappa_n] - \mathbb{E}_{z \sim \mathcal{D}_\textup{clean}}[\kappa_n]
      < \mathbb{E}_{z \sim \mathcal{D}_\textup{noisy}}[\kappa^\dagger_n] - \mathbb{E}_{z \sim \mathcal{D}_\textup{clean}}[\kappa^\dagger_n].
      \label{eq:prop2}
    \end{equation}
  \end{proposition}
  \begin{proof}[Proof sketch (full proof: Appendix~A.2)]
    Define $\Delta_t := \mathbb{E}_\textup{noisy}[1/I_{(n,t)}] - \mathbb{E}_\textup{clean}[1/I_{(n,t)}]$.
    The left-hand side of Eq.~\ref{eq:prop2} equals $\sum_t k(t)\Delta_t$.
    For $t \leq T_m$ (before memorization onset): clean gradients decay so $\Delta_t \leq 0$, strictly at some $t^*$ (Appendix~A.2).
    Beyond $T_m$: noisy gradients also drop (memorization), so $\Delta_{T-1} > 0$ — the last-checkpoint baseline is inflated by memorization.
    The decreasing kernel upweights $t \leq T_m$ (where $\Delta_t < 0$) and downweights $t > T_m$ (where $\Delta_t > 0$).
    Since $k$ is strictly decreasing, the negative-$\Delta_t$ early terms receive more weight and the large positive $\Delta_{T-1}$ term receives the least, making the weighted sum strictly smaller than $\Delta_{T-1} > 0$.
  \end{proof}
\end{foldable}

\subsection{Synthetic Candidate Pool Generation} \label{sec:04-method-pool}

\begin{foldable} % Why a synthetic pool
  Rather than selecting directly from $\mathcal{D}$, TAKE draws $\tilde{\mathcal{D}}$ from a synthetic candidate pool $\mathcal{D}_\text{synth}$.
  Direct subsampling has two problems: at extreme compression, verbatim records raise privacy concerns, and the budget $K$ is too small to cover the full training support.
  We generate $\mathcal{D}_\text{synth}$ by fine-tuning a small language model on $\mathcal{D}$ and sampling from it.
  The fine-tuned LM generalizes beyond memorized training instances, and stochastic sampling produces a diverse pool that extends beyond the training support.
  $\mathcal{D}_\text{synth}$ serves solely as the selection pool and is never appended to downstream training data.
\end{foldable}

\begin{foldable} % Distribution mismatch is beneficial
  TAKE does \emph{not} assume $P_{\mathcal{D}_\text{synth}} = P_{\mathcal{D}}$; this distributional gap is \emph{beneficial}.
  First, coverage expansion fills undersampled regions of $\mathcal{D}$, improving generalization of the distilled corpus.
  Second, model fluency bias suppresses noisy samples that HSB would otherwise over-select, yielding a cleaner candidate pool.
  Third, the distribution mismatch physically decouples $\tilde{\mathcal{D}}$ from original records, reducing verbatim exposure risk.
  OT handles the distributional gap gracefully: it routes mass toward $P^w$ regardless of pool distribution, so $P_{\tilde{\mathcal{D}}}$ approximates $P^w$ provided $\mathcal{D}_\text{synth}$ covers the support of $P_{\mathcal{D}}$ (Assumption~\ref{ass:coverage}):
  \begin{assumption}[Pool Coverage] \label{ass:coverage}
    For every training sample $z_n \in \mathcal{D}$, there exists a synthetic sample $\tilde{z}_c \in \mathcal{D}_\text{synth}$ within embedding distance $r$:
    $\min_c \|\phi(z_n) - \phi(\tilde{z}_c)\| \leq r$.
  \end{assumption}
  Under Assumption~\ref{ass:coverage}, $W_1(P^w, P_{\mathcal{D}_\text{synth}}) \leq r$, so the optimal transport cost is at most $r$ before any prototype selection.
  A practitioner can monitor coverage directly by reporting the mean nearest-neighbor distance $\bar{r} = \frac{1}{N}\sum_n \min_c \|\phi(z_n) - \phi(\tilde{z}_c)\|$.
\end{foldable}

\subsection{Distributional Prototype Selection via Discrete Optimal Transport} \label{sec:04-method-ot}

\begin{foldable} % Why discrete OT
  Both $\mathcal{D}$ and $\mathcal{D}_\text{synth}$ are finite sample sets, so the matching problem is naturally over empirical measures.
  Discrete OT operates directly on $\mu = \sum_n w_n \delta_{z_n}$ and $\nu = \frac{1}{|\mathcal{D}_\text{synth}|}\sum_c \delta_{\tilde{z}_c}$, requiring no density estimation and yielding an explicit transport plan with direct prototype-level assignments.
\end{foldable}

\begin{foldable} % Formulation and entropic regularization
  Let $C_{nc} = \|\phi(z_n) - \phi(\tilde{z}_c)\|^2$ be the cost matrix in a task-relevant embedding space.
  Bare Kantorovich OT, $\min_{\gamma \in \Pi(\mu,\nu)} \langle C, \gamma \rangle$, produces a sparse, non-differentiable plan and scales as $O(N^3)$.
  We add entropic regularization to obtain a strictly convex, smooth objective solvable via Sinkhorn--Knopp iterations at $O(N \cdot |\mathcal{D}_\text{synth}|)$ per step:
  \begin{equation}
    \gamma^* = \arg\min_{\gamma \in \Pi(\mu,\nu)} \langle C, \gamma \rangle + \varepsilon \cdot \mathrm{KL}(\gamma \| \mu \otimes \nu).
    \label{eq:ot}
  \end{equation}
  The soft coupling distributes mass across nearby candidates rather than hard one-to-one assignment, reducing sensitivity to embedding errors and encouraging diversity in $\tilde{\mathcal{D}}$.
  The $K$ candidates with highest received mass $\sum_n \gamma^*_{nc}$ are selected as $\tilde{\mathcal{D}}$, resolving open problem~\ref{open-problem-select}.
  The following theorem closes the theoretical chain from RDM (Eq.~\ref{eq:reweighted-distribution-matching}) back to DD (Eq.~\ref{eq:dataset-distillation}), bounding the downstream loss gap in terms of the OT cost.

  \begin{theorem}[Distribution Matching Gap Bound] \label{thm:gap}
    Let $f^*_{\mathcal{S}}$ denote the model trained on $\mathcal{S}$, and let $\mathcal{L}(\mathcal{S}) = \mathbb{E}_{z \sim P^w}[\ell(z;\,f^*_{\mathcal{S}})]$.
    Assume: \emph{(i)} $\ell(\cdot;\,f)$ is $L$-Lipschitz in the embedding metric for all $f$; \emph{(ii)} the model class has uniform stability constant $\beta_K = O(1/K)$~\cite{bousquet2002stability}, so that $\sup_z |\ell(z;\,f^*_{\tilde{\mathcal{D}}}) - \ell(z;\,f^*_\mathcal{D})| \leq \beta_K$.
    Then:
    \begin{equation}
      \bigl|\mathcal{L}(\tilde{\mathcal{D}}) - \mathcal{L}(\mathcal{D})\bigr|
      \;\leq\; 2L \cdot W_1(P^w,\,P_{\tilde{\mathcal{D}}}) + \beta_K,
      \label{eq:gap-bound}
    \end{equation}
    where $W_1$ is the Wasserstein-1 distance under the embedding metric.
    As $K \to \infty$, $\beta_K \to 0$, and the bound is dominated by the $W_1$ term controlled by the OT objective.
  \end{theorem}
  \begin{proof}[Proof sketch (full proof: Appendix~A.3)]
    Insert two intermediate terms, $\mathbb{E}_{P_{\tilde{\mathcal{D}}}}[\ell(z;\,f^*_{\tilde{\mathcal{D}}})]$ and $\mathbb{E}_{P_{\tilde{\mathcal{D}}}}[\ell(z;\,f^*_\mathcal{D})]$, and apply the triangle inequality twice.
    This yields two distribution-shift terms (under fixed $f^*_{\tilde{\mathcal{D}}}$ and $f^*_\mathcal{D}$ respectively), each bounded by $L \cdot W_1(P^w, P_{\tilde{\mathcal{D}}})$ via Kantorovich--Rubinstein duality; and one model-difference term under $P_{\tilde{\mathcal{D}}}$, bounded by $\beta_K = O(1/K)$ via uniform stability.
  \end{proof}
  The bound is stated with respect to $P^w$ rather than the uniform $P_\mathcal{D}$ --- this is intentional, as $P^w$ is the designed optimization target of the RDM objective; since $P^w$ reweights $P_\mathcal{D}$ over the same support, minimizing $W_1(P^w, P_{\tilde{\mathcal{D}}})$ controls coverage of $P_\mathcal{D}$ up to the reweighting gap bounded in Remark~\ref{rem:double-gap}.
\end{foldable}

\begin{foldable} % Method wrap-up
  Together, Propositions~\ref{prop:hsb}--\ref{prop:traj} and Theorem~\ref{thm:gap} establish the full theoretical basis for TAKE: HSB is structural and unavoidable at a single checkpoint (Proposition~\ref{prop:hsb}); trajectory integration strictly reduces it (Proposition~\ref{prop:traj}); and minimizing the OT cost directly controls the downstream loss gap under $P^w$ (Theorem~\ref{thm:gap}; the gap to $P_\mathcal{D}$ is bounded by Remark~\ref{rem:double-gap}).
  Algorithm~\ref{alg:take} consolidates the full TAKE pipeline.
\end{foldable}

\begin{center}
  \begin{minipage}{0.99\linewidth}
    \begin{algorithm}[H]
      \SetKwInOut{KwIn}{Input}
      \SetKwInOut{KwOut}{Output}
      \SetKwComment{Comment}{}{}

      \caption{Trajectory-Aware Knowledge Estimation (\textbf{TAKE})}
      \label{alg:take}

      \KwIn{corpus $\mathcal{D} = \{z_n\}_{n=1}^{N}$; budget $K$; checkpoints $T$; kernel $k(\cdot)$; $\varepsilon$; pool size $M$}
      \KwOut{distilled corpus $\tilde{\mathcal{D}}$, $|\tilde{\mathcal{D}}| = K$}

      \vspace{-0.5\baselineskip}
      \Comment{\noindent\hrulefill}

      \vspace{0.5\baselineskip}
      \Comment{$\triangleright$ \textbf{\textit{Stage 1: Score}}}
      Train logistic probe on $\mathcal{D}$, saving checkpoints $\theta_0, \ldots, \theta_{T-1}$ \\
      \For{$t = 0, \ldots, T-1$}{
        Compute $I_{(n,t)} \leftarrow \|\nabla_W \ell(z_n;\,\theta_t)\|$ for all $n$ \hfill\textit{(Goodfellow trick)} \\
        Normalize $I_{(\cdot,t)} \leftarrow I_{(\cdot,t)} / \sum_{n'} I_{(n',t)}$ \hfill\textit{(column-wise)}
      }
      Compute $\kappa_n \leftarrow \textstyle\sum_{t=0}^{T-1} k(t) / I_{(n,t)}$ for all $n$ \hfill via Eq.~\ref{eq:kappa} \\
      Normalize $w_n \leftarrow \kappa_n / \sum_{n'} \kappa_{n'}$ for all $n$ \hfill\textit{(defines $P^w$)}

      \vspace{0.5\baselineskip}
      \Comment{$\triangleright$ \textbf{\textit{Stage 2: Generate}}}
      Fine-tune small language model $\mathrm{LM}$ on $\mathcal{D}$ \\
      Sample synthetic pool $\mathcal{D}_\text{synth} = \{\tilde{z}_c\}_{c=1}^{M}$ from $\mathrm{LM}$

      \vspace{0.5\baselineskip}
      \Comment{$\triangleright$ \textbf{\textit{Stage 3: Distill}}}
      Embed: $\phi(z_n)$ for $n \in [N]$; $\phi(\tilde{z}_c)$ for $c \in [M]$ \\
      Compute cost matrix $C_{nc} = \|\phi(z_n) - \phi(\tilde{z}_c)\|^2$ \\
      Solve $\gamma^* \leftarrow \mathrm{Sinkhorn}(C,\, w,\, \mathbf{1}/M,\, \varepsilon)$ \hfill via Eq.~\ref{eq:ot} \\
      Distill $\tilde{\mathcal{D}} \leftarrow \mathrm{top}\text{-}K\!\left(\{\tilde{z}_c\},\; \textstyle\sum_n \gamma^*_{nc}\right)$ \hfill\textit{(by received mass)}

      \vspace{0.5\baselineskip}
      \Return{$\tilde{\mathcal{D}}$}
    \end{algorithm}
  \end{minipage}
\end{center}
